\documentclass[12pt,a4paper]{article}
\usepackage[utf8]{inputenc}
\usepackage[T2A]{fontenc}
\usepackage[russian]{babel}
\usepackage{amsmath,amssymb}
\usepackage{geometry}
\geometry{left=2.5cm,right=2.5cm,top=2.5cm,bottom=2.5cm}
\usepackage{graphicx}
\usepackage{booktabs}
\usepackage{hyperref}
\usepackage{listings}
\usepackage{xcolor}

\lstset{
    basicstyle=\ttfamily\small,
    keywordstyle=\color{blue},
    commentstyle=\color{gray},
    frame=single,
    breaklines=true,
    columns=fullflexible
}

\title{Отчет по практическому заданию №1}
\author{Гусейнов А.Э. 681}

\begin{document}

\maketitle

\tableofcontents
\newpage

\section{Введение}

В работе рассматривается двумерная задача Дирихле для уравнения Пуассона в криволинейной области. Для приближенного решения используется метод фиктивных областей с продолжением по младшему коэффициенту. Полученная дифференциальная задача аппроксимируется на равномерной прямоугольной сетке разностной схемой. Система линейных алгебраических уравнений решается методом сопряженных градиентов с диагональным предобуславливанием. Последовательная программа распараллелена с помощью технологии OpenMP.

Вариант задания: \textbf{вариант 7} — квадрат с отсеченной вершиной:
\[
D = \{(x,y): |x|+|y|<2,\; y<1\}.
\]
Вершины области: \((-2,0)\), \((0,-2)\), \((2,0)\), \((1,1)\), \((-1,1)\).

\section{Постановка задачи}

В области \(D\subset\mathbb{R}^2\) с кусочно-гладкой границей \(\gamma\) требуется найти функцию \(u(x,y)\), удовлетворяющую уравнению Пуассона
\begin{equation}
-\Delta u = f(x,y), \qquad (x,y)\in D,
\label{eq:poisson}
\end{equation}
и граничному условию Дирихле
\begin{equation}
u(x,y) = \varphi(x,y), \qquad (x,y)\in\gamma.
\label{eq:bc}
\end{equation}
Для варианта 7:
\[
f(x,y)=1,\qquad \varphi(x,y)=\sin(x+y).
\]

\section{Метод фиктивных областей}

Область \(D\) погружается в прямоугольник
\[
\Pi = \{(x,y): A_1<x<B_1,\; A_2<y<B_2\}.
\]
Для варианта 7: \(A_1=-2\), \(B_1=2\), \(A_2=-2\), \(B_2=1\). Фиктивная область \(\hat D = \Pi\setminus D\). Вводится малый параметр \(\varepsilon>0\). Решается задача
\begin{equation}
-\Delta v + Q(x,y)\bigl(v-v_0(x,y)\bigr) = F(x,y), \quad (x,y)\in\Pi\setminus\gamma,
\label{eq:fict}
\end{equation}
\[
v(x,y)=v_0(x,y),\quad (x,y)\in\Gamma,
\]
где
\[
Q(x,y)=
\begin{cases}
0, & (x,y)\in D,\\[2mm]
\varepsilon^{-2}, & (x,y)\in\hat D,
\end{cases}
\qquad
F(x,y)=
\begin{cases}
f(x,y), & (x,y)\in D,\\[2mm]
0, & (x,y)\in\hat D.
\end{cases}
\]
В качестве продолжения \(v_0\) используется \(\varphi(x,y)=\sin(x+y)\). Параметр \(\varepsilon\) выбирается как \(\varepsilon=h^2\), где \(h=\max(h_1,h_2)\) — максимальный шаг сетки.

\section{Разностная схема}

В прямоугольнике \(\Pi\) вводится равномерная сетка
\[
\bar\omega_h = \bar\omega_1\times\bar\omega_2,\qquad
\bar\omega_1=\{x_i=A_1+ih_1,\; i=0,\dots,M\},\quad
\bar\omega_2=\{y_j=A_2+jh_2,\; j=0,\dots,N\},
\]
где \(h_1=(B_1-A_1)/M\), \(h_2=(B_2-A_2)/N\). Внутренние узлы обозначаются \(\omega_h\). Дифференциальное уравнение аппроксимируется разностным уравнением
\begin{equation}
-w_{x\bar x,ij} - w_{y\bar y,ij} + Q_{ij}w_{ij}
= F_{ij} + Q_{ij}v_0(x_i,y_j),
\qquad i=1,\dots,M-1,\quad j=1,\dots,N-1.
\label{eq:scheme}
\end{equation}
Граничные условия аппроксимируются точно:
\[
w_{ij}=v_0(x_i,y_j),\qquad (x_i,y_j)\in\Gamma.
\]
Коэффициенты \(Q_{ij}\) и \(F_{ij}\) вычисляются как средние по ячейке \(\Pi_{ij}\):
\[
Q_{ij}=\frac{1}{h_1h_2}\iint_{\Pi_{ij}}Q(x,y)\,dx\,dy,\qquad
F_{ij}=\frac{1}{h_1h_2}\iint_{\Pi_{ij}}F(x,y)\,dx\,dy.
\]
Для варианта 7 доля площади ячейки, попавшая в \(D\), вычислялась приближенно по подсетке \(12\times12\), что позволяет единообразно обрабатывать криволинейные границы. В результате получается система линейных алгебраических уравнений
\begin{equation}
Aw=B
\label{eq:system}
\end{equation}
с симметричной положительно определенной матрицей \(A\).

\section{Метод сопряженных градиентов с диагональным предобуславливанием}

Система \eqref{eq:system} решается методом сопряженных градиентов. Используется диагональный предобуславливатель
\[
(Dw)_{ij}=\left(\frac{2}{h_1^2}+\frac{2}{h_2^2}+Q_{ij}\right)w_{ij}.
\]
В задании скалярное произведение и соответствующая ему евклидова норма определены как
\[
(u,v)_E=h_1h_2\sum_{i=1}^{M-1}\sum_{j=1}^{N-1}u_{ij}v_{ij},
\qquad
\|u\|_E=\sqrt{(u,u)_E}.
\]
Итерационный процесс:
\begin{align*}
&r^{(0)}=B-Aw^{(0)},\qquad Dz^{(0)}=r^{(0)},\qquad p^{(1)}=z^{(0)},\\
&\alpha_k=\frac{(z^{(k-1)},r^{(k-1)})}{(Ap^{(k)},p^{(k)})},\\
&w^{(k)}=w^{(k-1)}+\alpha_k p^{(k)},\\
&r^{(k)}=r^{(k-1)}-\alpha_k Ap^{(k)},\\
&Dz^{(k)}=r^{(k)},\\
&\beta_{k+1}=\frac{(z^{(k)},r^{(k)})}{(z^{(k-1)},r^{(k-1)})},\\
&p^{(k+1)}=z^{(k)}+\beta_{k+1}p^{(k)}.
\end{align*}
Начальное приближение \(w^{(0)}=0\). Критерий остановки:
\[
\|w^{(k+1)}-w^{(k)}\|_E<\delta,
\]
где \(\delta=10^{-8}\). В отношениях для \(\alpha_k\) и \(\beta_k\) общий множитель \(h_1h_2\) сокращается, поэтому в программе для них используются обычные суммы. Критерий остановки, напротив, проверяется в требуемой сеточной норме:
\[
\|w^{(k+1)}-w^{(k)}\|_E
=\sqrt{h_1h_2}\,|\alpha_k|\sqrt{\sum_{i,j}(p_{ij}^{(k)})^2}<\delta.
\]

Качество итераций дополнительно контролируется функционалом
\[
J(w)=\frac12(Aw,w)_E-(B,w)_E.
\]
Это числовая «энергия» сеточного решения, а не дополнительная неизвестная. Так как \(A\) положительно определен, минимум \(J\) достигается при \(Aw=B\). В программе функционал вычисляется непосредственно по формуле; если его рост превышает оценку погрешности суммирования, решение возвращается к последней итерации с монотонным убыванием, остаток пересчитывается, а метод сопряженных градиентов перезапускается с новым направлением.

\section{Параллельная реализация на OpenMP}

Последовательная и параллельная версии собраны из одного файла \texttt{main.cpp}: без флага \texttt{-fopenmp} и с этим флагом соответственно. Распараллелены вычисление \(Q_{ij}\), \(F_{ij}\), правой части \(B\), умножение матрицы на вектор, скалярные произведения и обновление векторов. Для вложенных циклов используется \texttt{collapse(2)}, распределение итераций — \texttt{schedule(static)}, для сумм — \texttt{reduction}. В OpenMP-версии время измеряется функцией \texttt{omp\_get\_wtime()} только внутри итерационного процесса, включая дополнительное умножение на \(A\) для проверки \(J\); построение сетки, коэффициентов и правой части в измеряемое время не входит.

Ключевой фрагмент умножения матрицы на вектор:
\begin{lstlisting}[language=C++]
#pragma omp parallel for collapse(2) schedule(static)
for (int i = 1; i <= M - 1; ++i) {
    for (int j = 1; j <= N - 1; ++j) {
        int k = idx(i, j);
        double val = (2.0*inv_h1_2 + 2.0*inv_h2_2 + Q[k]) * u[k];
        if (i > 1)     val -= u[idx(i-1,j)] * inv_h1_2;
        if (i < M - 1) val -= u[idx(i+1,j)] * inv_h1_2;
        if (j > 1)     val -= u[idx(i,j-1)] * inv_h2_2;
        if (j < N - 1) val -= u[idx(i,j+1)] * inv_h2_2;
        out[k] = val;
    }
}
\end{lstlisting}

Компиляция:
\begin{lstlisting}[language=bash]
g++ -O0 -std=c++17 main.cpp -o solve
g++ -O0 -fopenmp -std=c++17 main.cpp -o solve_omp
\end{lstlisting}
Эти команды соответствуют локальному скрипту сборки \texttt{build.bat}; параметры сборки исполняемых файлов, использованных для расчетов в кластере, в сохраненных логах не указаны.
Для сетки \(80\times80\) запуск выполнялся с 1, 4 и 16 нитями:
\begin{lstlisting}[language=bash]
export OMP_NUM_THREADS=1
./solve_omp 80 80
export OMP_NUM_THREADS=4
./solve_omp 80 80
export OMP_NUM_THREADS=16
./solve_omp 80 80
\end{lstlisting}

\section{Результаты расчетов}

В таблице~\ref{tab:seq} приведены последовательные расчеты, в таблицах~\ref{tab:omp80}--\ref{tab:omp800} — ранее сохраненные OpenMP-запуски из папки \texttt{result}. Эти файлы получены до исправления критерия остановки и добавления контроля \(J\), поэтому времена, итерации и невязки в таблицах не являются результатами исправленной версии и должны быть пересчитаны для итогового сравнения. Время — время итерационного процесса; невязка исправленной программы выводится в сеточной норме \(\|r\|_E\). Число неизвестных равно \((M-1)(N-1)\).

\begin{table}[h]
\centering
\caption{Архивные последовательные расчеты до исправления критерия остановки}
\label{tab:seq}
\begin{tabular}{crrr}
\toprule
\(M\times N\) & Итерации & Время, с & \(\|r\|_2\)\\
\midrule
\(20\times20\) & 9 & 0.000936245 & \(4.59617\cdot10^{-7}\)\\
\(40\times40\) & 7 & 0.00305351 & \(1.28974\cdot10^{-7}\)\\
\(80\times80\) & 5 & 0.0104305 & \(4.31385\cdot10^{-6}\)\\
\bottomrule
\end{tabular}
\end{table}

\begin{table}[h]
\centering
\caption{Архивная масштабируемость OpenMP на сетке \(80\times80\), до исправления}
\label{tab:omp80}
\begin{tabular}{crrr}
\toprule
Число нитей & Итерации & Время, с & Ускорение относительно 1 нити\\
\midrule
1  & 5 & 0.00896651 & 1.000\\
4  & 5 & 0.00235089 & 3.814\\
16 & 5 & 0.000963256 & 9.309\\
\bottomrule
\end{tabular}
\end{table}

\begin{table}[h]
\centering
\caption{Архивная масштабируемость OpenMP на сетке \(400\times600\), до исправления}
\label{tab:omp400}
\begin{tabular}{crrr}
\toprule
Число нитей & Итерации & Время, с & Ускорение относительно 2 нитей\\
\midrule
2  & 4 & 0.135180 & 1.000\\
4  & 4 & 0.0677342 & 1.996\\
8  & 4 & 0.0345109 & 3.917\\
16 & 4 & 0.0228812 & 5.908\\
\bottomrule
\end{tabular}
\end{table}

\begin{table}[h]
\centering
\caption{Архивная масштабируемость OpenMP на сетке \(800\times1200\), до исправления}
\label{tab:omp800}
\begin{tabular}{crrr}
\toprule
Число нитей & Итерации & Время, с & Ускорение относительно 4 нитей\\
\midrule
4  & 4 & 0.288312 & 1.000\\
8  & 4 & 0.135952 & 2.121\\
16 & 4 & 0.0587516 & 4.907\\
32 & 4 & 0.0950004 & 3.035\\
\bottomrule
\end{tabular}
\end{table}

В архивных запусках прежней версии программы на сетке \(80\times80\) ускорение относительно одного потока составило \(3.814\) для 4 потоков и \(9.309\) для 16 потоков. На сетке \(400\times600\) рост числа потоков с 2 до 16 дал ускорение \(5.908\), а на сетке \(800\times1200\) наилучшее время получено при 16 потоках. При 32 потоках время увеличилось относительно 16 потоков. Эти измерения характеризуют только прежнюю версию; после исправления критерия остановки и добавления контроля \(J\) производительность следует измерить заново. Время разных размеров сетки напрямую сравнивать нельзя.

В этих запусках максимальное число итераций ограничено числом неизвестных. Результаты OpenMP-запусков могут незначительно отличаться по невязке из-за порядка суммирования при параллельной редукции. Значения невязки характеризуют решение дискретной системы с указанным критерием остановки и не являются оценкой ошибки относительно точного решения исходной дифференциальной задачи.

\section{Заключение}

Реализованы метод фиктивных областей, конечно-разностная схема и решение системы методом сопряженных градиентов с диагональным предобуславливанием. Основные циклы распараллелены с помощью OpenMP. Исправленная программа использует норму из задания и контролирует монотонность функционала \(J\), при необходимости перезапуская итерационный процесс. Ранее сохраненные измерения относятся к предыдущей версии программы; их необходимо повторить, чтобы корректно сравнить производительность и невязку исправленной версии.

\begin{thebibliography}{9}
\bibitem{tihonov}
А.Н.~Тихонов, А.А.~Самарский. Уравнения математической физики. М.: Наука, 1977.
\bibitem{marchuk}
Г.И.~Марчук. Методы вычислительной математики. М.: Наука, 1989.
\bibitem{vabishevich}
П.Н.~Вабищевич. Метод фиктивных областей в задачах математической физики. 2017.
\bibitem{samarsky}
А.А.~Самарский. Теория разностных схем. М.: Наука, 1989.
\bibitem{samarsky_gulin}
А.А.~Самарский, А.В.~Гулин. Численные методы. М.: Наука, 1989.
\bibitem{andreev}
В.Б.~Андреев. Численные методы. М.: Макс Пресс, 2013.
\bibitem{antonov}
А.С.~Антонов. Технологии параллельного программирования MPI и OpenMP. М.: Изд-во МГУ, 2012.
\bibitem{openmp}
OpenMP Architecture Review Board. OpenMP Application Programming Interface. \url{https://www.openmp.org/}.
\end{thebibliography}

\appendix
\section{Основные фрагменты программы}

Полный код программы прилагается отдельно. Ниже приведен фрагмент, задающий геометрию варианта 7 и функцию проверки точки внутри области.

\begin{lstlisting}[language=C++]
const vector<Point> DOMAIN = {
    {-2.0,  0.0},
    { 0.0, -2.0},
    { 2.0,  0.0},
    { 1.0,  1.0},
    {-1.0,  1.0}
};

double f(double x, double y)   { return 1.0; }
double phi(double x, double y) { return sin(x + y); }
double v0(double x, double y)  { return sin(x + y); }

bool insideDomain(double x, double y) {
    bool inside = false;
    int n = (int)DOMAIN.size();
    for (int i = 0, j = n - 1; i < n; j = i++) {
        const Point& pi = DOMAIN[i];
        const Point& pj = DOMAIN[j];
        bool intersect = ((pi.y > y) != (pj.y > y)) &&
            (x < (pj.x - pi.x) * (y - pi.y) / (pj.y - pi.y) + pi.x);
        if (intersect) inside = !inside;
    }
    return inside;
}
\end{lstlisting}

\end{document}